%% kapa: exam=1 round=33 year=2022 period=2 form=A subject=accounting subject_ko=회계학 seq=18 number=58
%% source: Q-Net 감정평가사 기출문제 (cst00309) artlSeq=5209997 file=raw/1008/5209997/2022년도 제33회 감정평가사 1차시험 2교시(감정평가 관계법규, 회계학).pdf
%% conversion: kordoc 4.24.51 -> tools/gen1.py (auto); review=proofread-sonnet; flags=none
\begin{question}{58}
고객과의 계약에서 생기는 수익에 관한 설명으로 옳지 않은 것은?
\begin{choices}
  \item 고객과의 계약에서 약속한 대가에 변동금액이 포함된 경우 기업은 고객에게 약속한 재화나 용역을 이전하고 그 대가로 받을 권리를 갖게 될 금액을 추정한다.
  \item 고객이 재화나 용역의 대가를 선급하였고 그 재화나 용역의 이전 시점이 고객의 재량에 따라 결정된다면, 기업은 거래가격을 산정할 때 화폐의 시간가치가 미치는 영향을 고려하여 약속된 대가(금액)를 조정해야 한다.
  \item 적절한 진행률 측정방법에는 산출법과 투입법이 포함되며, 진행률 측정방법을 적용할 때 고객에게 통제를 이전하지 않은 재화나 용역은 진행률 측정에서 제외한다.
  \item 고객과의 계약체결 증분원가가 회수될 것으로 예상된다면 이를 자산으로 인식한다.
  \item 고객이 기업이 수행하는 대로 기업의 수행에서 제공하는 효익을 동시에 얻고 소비한다면, 기업은 재화나 용역에 대한 통제를 기간에 걸쳐 이전하는 것이므로 기간에 걸쳐 수익을 인식한다.
\end{choices}
\end{question}
\answer{2}
